\documentclass[11pt]{article}
\usepackage[letterpaper,margin=1in]{geometry}
\usepackage{fancyhdr}
\usepackage{amsfonts}
\usepackage{amsthm}
\usepackage{amsmath}
\usepackage{amssymb}
\usepackage{newpxtext}
\usepackage{eulerpx}
\usepackage{eucal}
\usepackage[normalem]{ulem}
\pagestyle{fancy}
\fancyhf{}
\renewcommand{\headrulewidth}{0pt}
\renewcommand{\footrulewidth}{0pt}

\widowpenalty=1000
\clubpenalty=1000 

\setlength{\parindent}{0pt}
\setlength{\parskip}{1em}
    

\begin{document}
\begin{center}
    Personal Statement
\end{center}

I am applying to Cal Poly's Master of Science in Mathematics program to strengthen my preparation for a Ph.D. in mathematics. As I near completion of my undergraduate degree at Occidental College, I see this program as an opportunity to demonstrate my readiness for advanced study, while also enhancing my research and teaching skills.

Although my interest in pure mathematics began relatively late in my undergraduate career, I have since enrolled in rigorous courses such as Abstract Algebra I, Abstract Algebra II, Advanced Linear Algebra, Real Analysis I, and Real Analysis II. These courses covered fundamental topics, including material from Dummit \& Foote up to Galois Theory, as well as the topology of the real numbers, normed vector spaces, abstract metric spaces, measure theory, and Lebesgue integration. 

I participated in a small research project last spring which combined areas of graph theory, linear algebra, and abstract algebra. Our research involved finding pairs of vectors $\bigl(\vec{r},\vec{d}\bigr) \in (\mathbb{F}_p[x])^n \times (\mathbb{F}_p[x])^n$ satisfying:
\begin{equation}
        \bigl(\text{diag}(\vec{d}) - \text{A}(G)\bigr)\cdot\vec{r} = \vec{0},
\end{equation}
where $G$ is a path with $n$ nodes, $\text{A}(G)$ is its adjacency matrix, and $\gcd(r_1,r_2,\dots,r_n) = 1$. During this experience, I relied heavily on SageMath \textemdash a comprehensive Python library \textemdash to efficiently analyze large data sets of solutions for Equation (1). Despite the incomplete results, I value this experience as it gave me insight into the research process which extended beyond my usual coursework.

My teaching experience began at Pasadena City College, where I tutored and served as a supplemental instructor for the mathematics department. Having held instruction hours for large groups of students, I believe this experience solidified my interest in teaching, and I am therefore excited about Cal Poly's learn-by-doing teaching opportunities, which will be immensely beneficial as I prepare for a Ph.D.

I am particularly drawn to Cal Poly's paid summer research initiatives and the flexibility to develop that work into a thesis. Equally appealing is the option to complete a thesis in lieu of an oral exam as the graduate program's culminating experience. These opportunities will help me explore diverse fields—such as differential geometry, algebraic geometry, or operator theory—as I determine what I would most like to pursue in a Ph.D. Building on my background in commutative algebra and my emerging interest in functional analysis, I believe Cal Poly's offerings will help me refine my research interests while improving my overall preparedness for doctoral study.

My success in advanced coursework at Occidental College, alongside my research involvement and teaching experience, has prepared me to thrive in Cal Poly's M.S. in Mathematics program. I look forward to collaborating with faculty, deepening my research interests, and developing my instructional skills\textemdash which will be an ideal foundation for my eventual transition into a Ph.D. program.

\end{document}